\section{Incidence, resolvers, and the twelve-line carrier}
\label{app:incidence-resolvers}

\subsection{Characteristic-zero incidence formula}

On the complete chart, a line has row matrix
\[
 \begin{pmatrix}1&0&a&b\\0&1&c&d\end{pmatrix}
\]
and points \((s,t,as+ct,bs+dt)\).  The frozen lexicographic shape writes
\[
 a(x)=-\frac{h_a(x)}{A_a},\qquad
 b(x)=-\frac{h_b(x)}{A_b},\qquad
 c(x)=-\frac{h_c(x)}{A_c},\qquad d(x)=x,
\]
where the three denominators are nonzero constants.  Here (x=d), the raw
primitive eliminant is \(\widetilde g(d)\), its leading coefficient is
\(c_g\), and the monic eliminant and integral orbit coordinate are
\(g=c_g^{-1}\widetilde g\) and \(\alpha_i=c_gd_i\), respectively, as in
\cref{eq:monic-eliminant,eq:alpha-scaling}.  Two chart lines at
\(x\ne y\) meet precisely when
\begin{equation}
 (a(y)-a(x))(d(y)-d(x))
 -(b(y)-b(x))(c(y)-c(x))=0.
\label{eq:two-line-determinant}
\end{equation}
For \(\nu\in\{a,b,c\}\), put
\[
 D_\nu(x,y)=\frac{h_\nu(y)-h_\nu(x)}{y-x}.
\]
After dividing \cref{eq:two-line-determinant} by \((y-x)^2\) and clearing
the constant denominators, the exact incidence carrier is
\begin{equation}
 J(x,y)=-D_aA_bA_c-D_bD_cA_a.
\label{eq:divided-J}
\end{equation}
The polynomial \(J\) remains defined on the diagonal.  Its gcd has no
diagonal root and satisfies \cref{eq:incidence-identities}.  At each good
specialization all \(27\) copies of \(H_x\) equal the monic modular gcd.
Specialization gives the upper degree bound \(10\), while the exact factor
gives the lower bound.

The sixers are enumerated as six-element independent sets in the meeting
graph.  For a sixer \(S\), its opposite is
\[
 S^\vee=\{\ell\notin S:\ell\text{ meets exactly five lines in }S\}.
\]
Every exact sixer has \(|S^\vee|=6\), and pairing \(S\) with \(S^\vee\)
produces the \(36\) double-sixes in \cref{eq:configuration-counts}.

\subsection{Resolver reconstruction and irreducibility}

Resolver arrays run from constant to leading term and are monic.  Every CRT
prime is proved prime; the eliminant remains squarefree, denominators are
units, and the modular orbit product is multiplied back.  Signed
reconstruction requires a modulus exceeding twice the coefficient bound.

At a squarefree good reduction, the degree of any rational factor is a
proper subset sum of the modular factor degrees.  The two sets arising from
\cref{eq:resolver-factor-patterns} are disjoint, proving irreducibility of
both resolvers.  Separately, the exact configuration action binds their roots
all-and-only to the \(36\) double-sixes.

\begin{table}[t]
\centering
\small
\begin{tabularx}{\textwidth}{@{}l>{\raggedright\arraybackslash}X
 >{\raggedright\arraybackslash}X@{}}
\toprule
object & reconstruction witness & independent identity\\
\midrule
incidence \(H_x,Q_x\)
& characteristic-zero monic degrees \(10,17\); good primes
\(7,37,10^{50}+12477\)
& \(g=H_xQ_x\), \(H_x\mid J\), no diagonal, all modular graphs agree\\
\(R_\theta\)
& \(99\) primes; \(4951\)-digit CRT modulus
& orbit product, exact factors, incompatible subset sums\\
\(R_\delta\)
& \(198\) primes; \(9901\)-digit modulus exceeding the \(9858\)-digit
height requirement
& orbit product, exact factors, incompatible subset sums\\
\(\Acarrier\)
& \(1048\) primes; \(100609\)-digit modulus; \(13\times36\) coefficient
table
& \(28\) zero division remainders and \(28\) matching forward coefficients\\
\bottomrule
\end{tabularx}
\caption{Compact reconstruction ledger.  Large arrays are kept in exact
electronic witnesses rather than printed as decimal tables.}
\label{tab:resolver-ledger}
\end{table}

\subsection{Carrier identity over
\texorpdfstring{\(\Q(\theta_D)\)}{Q(theta-D)}}

Write
\[
 \Acarrier(d)=d^{12}+a_{11}d^{11}+\cdots+a_0,
 \qquad a_i\in\Q(\theta_D).
\]
The subtop coefficient has the exact normalization
\begin{equation}
 [d^{11}]\Acarrier=-\sum_{i=1}^{12}d_i=-\theta_D/c_g.
\label{eq:appendix-carrier-subtop}
\end{equation}
Exact long
division by \(\Acarrier\) yields \(\Bcarrier\) with zero remainder, and an
independent convolution recomputes all \(28\) product coefficients.  The
\([12,15]\) orbit partition binds the roots of \(\Acarrier\) to \(D\), not
merely to an arbitrary degree-\(12\) factor.

The stable carrier-table digest is
\[
\code{72d4aef5120926ec09904b08219cba7cf2b49323bd085d05274e5c17e1ed90a1}.
\]
This digest identifies the compact table but does not replace the division,
forward-product, or orbit-binding checks.
